\documentclass[11pt,a4paper]{article}

\usepackage[utf8]{inputenc}
\usepackage[T1]{fontenc}
\usepackage[french]{babel}
\usepackage{lmodern}
\usepackage{geometry}
\usepackage{microtype}
\usepackage{amsmath}
\usepackage{amssymb}
\usepackage{array}
\usepackage{booktabs}
\usepackage{longtable}
\usepackage{enumitem}
\usepackage{xcolor}
\usepackage{hyperref}

\geometry{margin=2.5cm}
\setlength{\parskip}{0.55em}
\setlength{\parindent}{0pt}
\setlist[itemize]{leftmargin=1.4em}
\setlist[enumerate]{leftmargin=1.6em}
\setlength{\emergencystretch}{3em}

\newcommand{\PC}{\mathrm{PC}}
\newcommand{\TPC}{\mathrm{TPC}}
\newcommand{\CC}{\mathrm{CC}}
\newcommand{\projectname}{\texttt{cc\allowbreak-formation\allowbreak-optimizer}}

\hypersetup{
  colorlinks=true,
  linkcolor=black,
  urlcolor=blue,
  pdftitle={Modélisation algorithmique complète - Optimisation des formations CC EAR 2027},
  pdfauthor={Projet cc-formation-optimizer}
}

\title{\textbf{Modélisation algorithmique complète}\\
\large Optimisation des formations de coordonnateurs communaux - EAR 2027}
\author{Projet \projectname}
\date{\today}

\begin{document}

\maketitle

\begin{abstract}
Ce document présente la modélisation algorithmique mise en oeuvre dans le
projet \projectname. Il décrit le problème métier, les
données d'entrée, les notations mathématiques, la formulation CP-SAT, le
pipeline de résolution, les contrôles de validation et les exports produits.
Il distingue explicitement les contraintes effectivement codées des alertes
d'interprétation et des évolutions métier possibles.
\end{abstract}

\tableofcontents
\newpage

\section{Introduction}

Le projet vise à organiser les sessions de formation des coordonnateurs
communaux pour les communes PC et TPC de l'EAR 2027. Chaque commune doit être
rattachée à une session, elle-même portée par une commune pivot. Une solution
acceptable doit respecter des contraintes opérationnelles fortes : capacité en
nombre de coordonnateurs communaux, temps de trajet maximal, budgets de
formations PC et TPC, cohérence de type entre communes PC et sessions TPC, et
existence effective des données de trajet.

Ce problème relève naturellement de l'optimisation sous contraintes. Il ne
s'agit pas seulement de construire des groupes arbitraires de communes : il
faut choisir quelles sessions ouvrir, décider à quel pivot rattacher chaque
commune, satisfaire les contraintes dures et arbitrer entre plusieurs solutions
faisables selon un objectif pondéré. Le modèle actuel est formulé comme un
programme linéaire en variables entières et résolu avec OR-Tools CP-SAT.

L'objectif du projet n'est pas uniquement mathématique. Le solveur fournit une
solution candidate, mais l'outil doit aussi produire une solution exploitable,
lisible et contrôlable par les acteurs métier. C'est pourquoi la chaîne
logicielle inclut une extraction métier, une validation indépendante, des
exports détaillés, des indicateurs d'alerte et une carte HTML de contrôle.

\section{Description métier du problème}

\subsection{Communes, catégories et coordonnateurs}

L'unité de base du problème est la commune. Chaque commune possède un code
INSEE, un nom, une population, une catégorie métier et éventuellement un
territoire EAR, un nombre de logements et des coordonnées géographiques. Les
catégories utilisées par le modèle sont PC et TPC.

Le modèle ne raisonne pas seulement en nombre de communes, mais en charge de
formation. Cette charge est exprimée en nombre de coordonnateurs communaux à
former. Une petite commune représente un coordonnateur, tandis qu'une commune
dont la population dépasse le seuil configuré représente deux coordonnateurs.
Cette règle donne plus de poids aux communes les plus peuplées dans les
contraintes de capacité et dans le coût de trajet.

\subsection{Sessions et pivots}

Une session de formation est identifiée par une commune pivot et par un rang de
slot. Toute commune chargée dans les données peut devenir pivot potentiel. Le
nombre de slots dépend de la catégorie de la commune pivot : une commune PC
peut porter jusqu'à trois sessions candidates, tandis qu'une commune TPC ne
peut porter qu'une session candidate.

Une session candidate ne devient réelle que si la variable d'ouverture
correspondante vaut 1. Le modèle actuel ne force pas une commune pivot ouverte
à être affectée à la session qu'elle porte. Ce choix est important : il permet
au solveur d'utiliser une commune comme lieu ou point d'organisation sans
imposer que cette commune fasse partie de ce groupe, mais il peut produire des
situations qui appellent un arbitrage métier.

\subsection{Territoires EAR et lisibilité métier}

Les territoires EAR sont chargés dans les données communales et apparaissent
dans les exports. Ils servent à qualifier les groupes, à calculer un territoire
majoritaire et à signaler des sessions multi-territoires ou des affectations à
un pivot d'un autre territoire. En revanche, dans l'état actuel du code, le
territoire EAR n'est pas une contrainte dure du modèle CP-SAT et n'intervient
pas directement dans la fonction objectif.

\subsection{Capacité, trajets et cohérence PC/TPC}

Chaque session ouverte doit contenir une charge comprise entre un minimum
\(L\) et une capacité maximale \(Q\), exprimée en nombre de CC. Les trajets sont
orientés : le temps de \(i\) vers \(j\) peut différer du temps de \(j\) vers
\(i\). Un couple commune-pivot est admissible seulement si le trajet existe
dans les données préparées et si son temps est inférieur ou égal au seuil
maximal \(T\).

La cohérence PC/TPC est gérée par une contrainte dure : une commune PC ne peut
pas être affectée à une session de type TPC. L'inverse n'est pas interdit :
une commune TPC peut être affectée à une session PC, mais cette mixité est
pénalisée dans l'objectif.

\section{Données d'entrée}

\subsection{Données communales}

Le solveur ne lit pas directement les fichiers bruts. La commande
\texttt{prepare-data} produit des CSV propres à partir des fichiers ODS ou CSV
du dossier brut. Le fichier communal propre contient notamment :

\begin{itemize}
  \item le code commune ;
  \item le nom de la commune ;
  \item la catégorie PC ou TPC ;
  \item le territoire EAR 2027 ;
  \item la population ;
  \item le nombre de logements ;
  \item la latitude et la longitude si elles sont disponibles.
\end{itemize}

La règle de calcul du nombre de CC est :
\begin{equation}
q_i =
\begin{cases}
1, & \text{si } p_i \leq p^{\star},\\
2, & \text{si } p_i > p^{\star},
\end{cases}
\label{eq:cc-count}
\end{equation}
où \(p_i\) désigne la population de la commune \(i\) et où le seuil configuré
est \(p^{\star}=5000\).

\subsection{Données de trajet}

Les temps de trajet sont préparés dans un CSV contenant une origine, un pivot
candidat et un temps en minutes. La préparation conserve l'orientation des
trajets et ne complète pas les valeurs absentes. Cette règle est structurante :
un trajet absent ne correspond pas à un coût très élevé, mais à une liaison
interdite car aucune variable d'affectation ne sera créée pour ce couple.

\subsection{Paramètres de configuration}

Les paramètres métier sont centralisés dans
\texttt{config/config\_ear2027.yaml}. Les valeurs courantes les plus
importantes sont :

\begin{center}
\begin{tabular}{lll}
\toprule
Paramètre & Valeur courante & Rôle \\
\midrule
\(T\) & 75 & temps maximal admissible en minutes \\
\(Q\) & 14 & capacité maximale d'une session en CC \\
\(L\) & 6 & remplissage minimal d'une session ouverte \\
\(B\) & 55 & budget total de sessions \\
\(f\) & 45 & budget de sessions PC \\
\(k\) & 10 & budget de sessions TPC \\
\(w_t\) & 1 & poids du trajet \\
\(w_e\) & 1000 & poids des coûts d'éligibilité \\
\(w_m\) & 500 & poids de la mixité résiduelle \\
\bottomrule
\end{tabular}
\end{center}

La configuration impose aussi \(B=f+k\), \(M_{\PC}=3\), \(M_{\TPC}=1\) et la
règle de population de l'équation~\eqref{eq:cc-count}.

\subsection{Données dérivées}

À partir des CSV propres et du YAML, le code construit les ensembles, les
variables candidates, les admissibilités, les compatibilités par défaut, les
coûts d'éligibilité et le nombre de slots par pivot. Ces objets dérivés sont
créés dans \texttt{parameters.py} et transmis au constructeur du modèle.

\section{Notations mathématiques}

\subsection{Ensembles}

\begin{longtable}{>{\raggedright\arraybackslash}p{0.18\textwidth}p{0.72\textwidth}}
\toprule
Notation & Signification \\
\midrule
\endhead
\(C\) & ensemble des communes à affecter. \\
\(P \subset C\) & ensemble des communes de catégorie PC. \\
\(R \subset C\) & ensemble des communes de catégorie TPC. La lettre \(R\) est utilisée ici pour éviter la confusion avec le seuil de trajet \(T\). \\
\(F=C\) & ensemble des communes candidates pivot. \\
\(M_j\) & nombre de slots candidats du pivot \(j\), avec \(M_j=3\) si \(j\in P\) et \(M_j=1\) si \(j\in R\). \\
\(S\) & ensemble des sessions candidates \((j,m)\), où \(j\in F\) et \(m\in\{1,\dots,M_j\}\). \\
\(S_i\) & ensemble des sessions candidates pour lesquelles une variable d'affectation de la commune \(i\) existe effectivement. \\
\bottomrule
\end{longtable}

\subsection{Paramètres}

\begin{longtable}{>{\raggedright\arraybackslash}p{0.18\textwidth}p{0.72\textwidth}}
\toprule
Notation & Signification \\
\midrule
\endhead
\(p_i\) & population de la commune \(i\). \\
\(q_i\) & nombre de CC associés à la commune \(i\). \\
\(\tau_{ij}\) & temps de trajet orienté de la commune \(i\) vers le pivot \(j\). \\
\(T\) & seuil maximal admissible de trajet. \\
\(a_{ij}\) & indicateur d'admissibilité du trajet \(i\to j\). \\
\(b_{ij}\) & indicateur de compatibilité métier du couple \(i,j\). \\
\(Q\) & capacité maximale d'une session ouverte. \\
\(L\) & remplissage minimal d'une session ouverte. \\
\(B\) & budget total déclaré de sessions. \\
\(f\) & budget maximal de sessions PC. \\
\(k\) & budget maximal de sessions TPC. \\
\(e_j^{\PC}\) & coût d'éligibilité du pivot \(j\) pour une session PC. \\
\(e_j^{\TPC}\) & coût d'éligibilité du pivot \(j\) pour une session TPC. \\
\(w_t,w_e,w_m\) & poids respectifs du trajet, de l'éligibilité et de la mixité. \\
\bottomrule
\end{longtable}

L'admissibilité et la compatibilité sont définies par :
\begin{align}
a_{ij}
&=
\begin{cases}
1, & \text{si le trajet } i\to j \text{ existe et } \tau_{ij}\leq T,\\
0, & \text{sinon,}
\end{cases}
\label{eq:admissibility}\\
b_{ij}
&\in \{0,1\}.
\label{eq:compatibility}
\end{align}
En l'absence de fichier de compatibilités, le code pose \(b_{ij}=1\) pour tous
les couples \(i,j\).

\subsection{Variables de décision}

Les variables effectivement créées sont les suivantes :

\begin{longtable}{>{\raggedright\arraybackslash}p{0.18\textwidth}p{0.72\textwidth}}
\toprule
Variable & Signification \\
\midrule
\endhead
\(x_{ijm}\in\{0,1\}\) & vaut 1 si la commune \(i\) est affectée à la session \((j,m)\). Cette variable n'existe que si \(a_{ij}=1\) et \(b_{ij}=1\). \\
\(y_{jm}\in\{0,1\}\) & vaut 1 si la session candidate \((j,m)\) est ouverte. \\
\(z_{jm}\in\{0,1\}\) & vaut 1 si la session ouverte \((j,m)\) est de type TPC ; si \(y_{jm}=1\) et \(z_{jm}=0\), elle est de type PC. \\
\(d_{jm}\in\{0,\dots,Q\}\) & variable entière représentant la mixité résiduelle, c'est-à-dire la charge TPC dans une session PC. \\
\bottomrule
\end{longtable}

\section{Formulation mathématique}

\subsection{Formulation synthétique}

Le modèle minimise une somme pondérée de trois composantes :
\begin{align}
\min \quad
& w_t O_{\mathrm{trajet}}
 + w_e O_{\mathrm{eligibilite}}
 + w_m O_{\mathrm{mixite}}
\label{eq:objective}
\end{align}
avec :
\begin{align}
O_{\mathrm{trajet}}
&=
\sum_{i\in C}
\sum_{(j,m)\in S_i}
q_i\,\tau_{ij}\,x_{ijm},
\label{eq:obj-travel}\\
O_{\mathrm{eligibilite}}
&=
\sum_{(j,m)\in S}
\left[
e_j^{\PC}y_{jm}
+
\left(e_j^{\TPC}-e_j^{\PC}\right)z_{jm}
\right],
\label{eq:obj-eligibility}\\
O_{\mathrm{mixite}}
&=
\sum_{(j,m)\in S} d_{jm}.
\label{eq:obj-mixing}
\end{align}

Les contraintes du modèle actuel sont :
\begin{align}
\sum_{(j,m)\in S_i} x_{ijm}
&=1
&& \forall i\in C,
\label{eq:unique-assignment}\\
x_{ijm}
&\leq y_{jm}
&& \forall i\in C,\ \forall (j,m)\in S_i,
\label{eq:opening}\\
L\,y_{jm}
&\leq
\sum_{i\in C:(j,m)\in S_i} q_i x_{ijm}
&& \forall (j,m)\in S,
\label{eq:min-fill}\\
\sum_{i\in C:(j,m)\in S_i} q_i x_{ijm}
&\leq
Q\,y_{jm}
&& \forall (j,m)\in S,
\label{eq:max-capacity}\\
z_{jm}
&\leq y_{jm}
&& \forall (j,m)\in S,
\label{eq:type-opening}\\
\sum_{(j,m)\in S} \left(y_{jm}-z_{jm}\right)
&\leq f,
\label{eq:budget-pc}\\
\sum_{(j,m)\in S} z_{jm}
&\leq k,
\label{eq:budget-tpc}\\
y_{j,m+1}
&\leq y_{jm}
&& \forall j\in P,\ \forall m\in\{1,\dots,M_j-1\},
\label{eq:slot-order}\\
x_{ijm}
&\leq 1-z_{jm}
&& \forall i\in P,\ \forall (j,m)\in S_i,
\label{eq:pc-not-in-tpc}\\
d_{jm}
&\geq
\sum_{i\in R:(j,m)\in S_i} q_i x_{ijm}
-Qz_{jm}
&& \forall (j,m)\in S,
\label{eq:mixing}\\
x_{ijm},y_{jm},z_{jm}
&\in\{0,1\},
\qquad
d_{jm}\in\{0,\dots,Q\}.
\label{eq:domains}
\end{align}

Le budget total \(B\) est vérifié en configuration par l'invariant \(B=f+k\).
Le modèle CP-SAT ne contient pas une troisième contrainte séparée
\(\sum y_{jm}\leq B\), car les contraintes~\eqref{eq:budget-pc} et
\eqref{eq:budget-tpc} pilotent directement les deux composantes du budget.

\subsection{Lecture des contraintes}

La contrainte~\eqref{eq:unique-assignment} impose que chaque commune soit
affectée exactement une fois. Si une commune ne dispose d'aucune session
admissible, la somme est vide et le modèle devient infaisable, ce qui reflète
un problème de données ou de paramétrage.

La contrainte~\eqref{eq:opening} interdit d'affecter une commune à une session
fermée. Les contraintes~\eqref{eq:min-fill} et~\eqref{eq:max-capacity}
encadrent la charge en CC de chaque session ouverte. Lorsqu'une session est
fermée, \(y_{jm}=0\) force sa charge à zéro.

La contrainte~\eqref{eq:type-opening} garantit qu'une session fermée ne peut
pas être déclarée TPC. Les contraintes~\eqref{eq:budget-pc} et
\eqref{eq:budget-tpc} limitent respectivement le nombre de sessions PC et TPC.
L'expression \(y_{jm}-z_{jm}\) vaut 1 pour une session PC ouverte, et 0 pour
une session TPC ou fermée.

La contrainte~\eqref{eq:slot-order} réduit les symétries pour les pivots PC :
un slot de rang supérieur ne peut être ouvert que si les slots précédents le
sont déjà. Cette règle ne change pas la nature des groupes possibles, mais elle
stabilise et simplifie la recherche du solveur.

La contrainte~\eqref{eq:pc-not-in-tpc} interdit les communes PC dans les
sessions TPC. Elle porte sur la catégorie des communes affectées, pas sur la
catégorie du pivot. Une session TPC peut donc être portée par un pivot PC si ce
pivot n'est pas affecté à cette session.

La contrainte~\eqref{eq:mixing} définit la variable de mixité. Si la session
est TPC, \(z_{jm}=1\), le terme \(-Qz_{jm}\) rend la borne non contraignante
grâce à la capacité maximale. Si la session est PC, \(z_{jm}=0\), et \(d_{jm}\)
couvre la charge TPC présente dans cette session PC.

\section{Fonction objectif}

La fonction objectif~\eqref{eq:objective} agrège trois critères réellement
présents dans le code.

Le terme de trajet~\eqref{eq:obj-travel} mesure le coût total des déplacements.
Chaque temps est pondéré par \(q_i\), de sorte qu'une commune représentant deux
coordonnateurs pèse deux fois plus qu'une commune représentant un seul
coordonnateur.

Le terme d'éligibilité~\eqref{eq:obj-eligibility} favorise les pivots jugés
plus adaptés selon les bandes de population configurées. La formulation est
linéaire : une session PC paie \(e_j^{\PC}\), tandis qu'une session TPC paie
\(e_j^{\TPC}\). Les coûts très élevés sont des pénalités, pas des interdictions
dures. Un pivot reste mathématiquement possible si les contraintes l'autorisent,
mais il devient très défavorable.

Le terme de mixité~\eqref{eq:obj-mixing} pénalise les CC TPC affectés à des
sessions PC. Le modèle autorise cette situation, car elle peut être nécessaire
pour respecter les capacités, les budgets et les trajets. Elle est toutefois
dégradée dans l'objectif afin d'être utilisée seulement si elle améliore
l'équilibre global.

Le modèle actuel ne contient pas de coût fixe d'ouverture de session, de
pénalité territoriale dure ou douce, de pénalité explicite pour pivot absent de
sa session, ni de contrainte de superviseur, de calendrier ou de disponibilité.
Ces éléments peuvent être étudiés comme extensions, mais ne doivent pas être
présentés comme des composantes actuelles.

\section{Méthode algorithmique}

\subsection{Rôle de CP-SAT}

Le solveur utilisé est OR-Tools CP-SAT. Il manipule des variables booléennes et
entières, des contraintes linéaires et une fonction objectif linéaire. Le modèle
ne contient pas de produit entre variables : les expressions de type, de budget
et d'éligibilité sont linéarisées directement.

La configuration transmet au solveur une limite de temps, un nombre de workers,
une graine aléatoire et l'activation éventuelle des logs. Ces paramètres
influencent la recherche, mais ne modifient pas le modèle mathématique.

\subsection{Pipeline de calcul}

La chaîne complète est volontairement séparée en étapes :

\begin{enumerate}
  \item chargement de la configuration YAML ;
  \item préparation des données brutes en CSV propres ;
  \item chargement des communes, trajets et compatibilités ;
  \item construction des paramètres dérivés ;
  \item diagnostic pré-résolution ;
  \item construction du modèle CP-SAT ;
  \item résolution stricte avec \texttt{solve}, ou résolution assouplie avec \texttt{solve-relaxed} ;
  \item extraction de la solution métier ;
  \item validation indépendante ;
  \item exports et carte HTML si demandés.
\end{enumerate}

Le workflow \texttt{solve-relaxed} commence par la configuration initiale, puis
teste une hiérarchie d'assouplissements explicitement configurés : modification
du poids de mixité, réduction des coûts TPC, augmentation du seuil de trajet,
réduction de \(L\), augmentation de \(Q\), augmentation cohérente de \(f\),
\(k\) et \(B\), puis remplacement final de coûts très élevés par une pénalité
finie. La contrainte PC vers session TPC n'est jamais relâchée automatiquement.

\subsection{Faisabilité, optimalité et validation}

Un statut \texttt{OPTIMAL} signifie que le solveur a trouvé une solution et
prouvé qu'aucune meilleure solution n'existe pour le modèle donné. Un statut
\texttt{FEASIBLE} signifie qu'une solution respectant les contraintes a été
trouvée, mais que la preuve d'optimalité n'est pas acquise dans les limites de
recherche.

Dans un contexte métier, une solution \texttt{FEASIBLE} peut être exploitable
si elle passe la validation et si ses indicateurs sont acceptables. Elle ne doit
cependant pas être décrite comme optimale. La validation métier reste distincte
de la validation algorithmique : le code vérifie les contraintes codées, mais
ne décide pas à la place des experts si une exception territoriale, un pivot
absent ou une mixité résiduelle est acceptable.

\section{Validation de la solution}

Le solveur retourne des valeurs de variables. Le module
\texttt{solution\_extractor.py} reconstruit ensuite des sessions ouvertes, des
affectations de communes et les composantes d'objectif. Cette solution extraite
est contrôlée par \texttt{validation.py} avant tout export.

Les contrôles effectivement réalisés sont :

\begin{itemize}
  \item chaque commune du modèle est affectée exactement une fois ;
  \item aucune affectation ne référence une session fermée ou inconnue ;
  \item chaque session respecte le remplissage minimal et la capacité maximale ;
  \item les budgets PC, TPC et total implicite sont respectés ;
  \item aucune commune PC n'est affectée à une session TPC ;
  \item chaque trajet utilisé existe et respecte le seuil \(T\) ;
  \item aucune compatibilité interdite n'est utilisée ;
  \item les types de session sont valides ;
  \item la variable de mixité est cohérente avec les affectations ;
  \item les composantes d'objectif et l'objectif total sont recalculés.
\end{itemize}

Les incohérences territoriales ne sont pas des violations bloquantes dans le
code actuel. Elles sont signalées dans les exports comme alertes métier. De
même, l'absence du pivot dans sa propre session n'est pas une violation
algorithmique aujourd'hui.

\section{Exports et interprétation}

\subsection{Fichiers produits}

Lorsque la validation réussit, le pipeline peut produire :

\begin{itemize}
  \item \texttt{sessions.csv}, avec une ligne par session ouverte ;
  \item \texttt{communes\_affectees.csv}, avec une ligne par commune affectée ;
  \item \texttt{rapport\_solution.md}, rapport lisible par un utilisateur métier ;
  \item \texttt{statistiques\_solution.json}, synthèse structurée ;
  \item \texttt{config\_utilisee.yaml}, copie de la configuration de résolution ;
  \item \texttt{solution\_formations.xlsx}, si \texttt{openpyxl} est disponible ;
  \item \texttt{solution\_map.html}, carte HTML autonome ;
  \item les journaux d'assouplissement si \texttt{solve-relaxed} est utilisé.
\end{itemize}

Les exports sont des produits de restitution. Ils ne remplacent pas la
validation initiale : ils sont écrits seulement après que la solution extraite a
été acceptée par \texttt{validate\_solution()}.

\subsection{Indicateurs de qualité}

Les exports permettent de suivre plusieurs indicateurs utiles :

\begin{itemize}
  \item nombre de sessions ouvertes et saturation des budgets ;
  \item nombre de sessions PC et TPC ;
  \item taux de remplissage et places restantes ;
  \item temps moyen et maximal ;
  \item sessions proches de la capacité maximale ;
  \item sessions faiblement remplies ;
  \item sessions multi-territoires ;
  \item mixité TPC dans les sessions PC ;
  \item pivots avec coût d'éligibilité élevé ;
  \item communes affectées à un pivot d'un autre territoire.
\end{itemize}

Ces indicateurs ne sont pas tous des contraintes. Ils servent à prioriser la
revue humaine et à identifier les arbitrages possibles.

\subsection{Carte HTML}

La carte HTML est un outil de contrôle visuel. Elle embarque les statistiques
globales, les contrôles de validation, les points cartographiés et une synthèse
des sessions. Elle utilise les coordonnées latitude/longitude des communes et
un fond IGN/Géoplateforme WMTS, tout en conservant les points visibles si le
fond externe ne charge pas.

La commande \texttt{render-map} permet de régénérer la carte depuis des exports
existants sans relancer le solveur. Cette distinction est importante lorsque la
résolution réelle est longue.

\section{Limites et points métier à arbitrer}

Le modèle actuel est volontairement centré sur l'affectation des communes aux
sessions. Plusieurs questions restent des choix métier, pas des erreurs du
solveur.

\subsection{Le pivot doit-il appartenir à sa propre session ?}

Aujourd'hui, aucune contrainte n'impose qu'une commune pivot soit affectée à la
session qu'elle porte. Imposer cette règle nécessiterait de garantir l'existence
des trajets diagonaux \(j\to j\) et d'ajouter une contrainte explicite. Selon le
niveau de rigidité souhaité, cette règle pourrait être une contrainte dure ou
une pénalité douce.

\subsection{Accepte-t-on les sessions multi-territoires ?}

Les territoires EAR sont visibles dans les exports, mais ne contraignent pas le
solveur. Si la lisibilité territoriale devient prioritaire, il faudrait ajouter
une pénalité ou une contrainte dédiée. Ce choix peut augmenter le coût de trajet
ou rendre certaines zones plus difficiles à affecter.

\subsection{Accepte-t-on des TPC dans des sessions PC ?}

Les TPC dans des sessions PC sont autorisées et pénalisées. À l'inverse, les PC
dans des sessions TPC sont interdites. Si la mixité résiduelle est jugée trop
élevée, il est possible de calibrer \(w_m\), d'ajuster les budgets TPC ou de
tester une contrainte plus stricte.

\subsection{Faut-il pénaliser davantage les trajets longs ?}

Le modèle interdit les trajets au-delà de \(T\) et minimise la somme pondérée
des temps admissibles. Il ne contient pas de pénalité spécifique pour les
trajets proches du seuil, en dehors de leur contribution linéaire au coût de
trajet. Une pénalité par palier pourrait être envisagée si les trajets proches
de \(T\) sont jugés trop nombreux.

\subsection{Faut-il viser moins de 55 sessions ?}

Le modèle respecte les plafonds \(f\), \(k\) et \(B\), mais ne pénalise pas
l'ouverture d'une session en tant que telle. Avec les paramètres actuels, une
solution peut donc saturer le plafond de 55 sessions si cela améliore les autres
termes de l'objectif. Si l'objectif métier est de réduire le nombre de sessions,
il faut ajouter un coût d'ouverture ou revoir les paramètres de capacité et de
remplissage.

\section{Conclusion}

Le modèle fournit une base reproductible et traçable pour organiser les
formations des coordonnateurs communaux. Il transforme des données brutes en
tables propres, construit une formulation CP-SAT cohérente avec les règles
métier codées, produit une solution candidate, puis la valide et l'exporte dans
des formats exploitables.

Son intérêt principal est de rendre explicites les arbitrages : temps de trajet,
capacité, budgets, éligibilité des pivots et mixité PC/TPC. L'outil ne remplace
pas l'arbitrage métier. Il fournit une solution contrôlée, documentée et
visualisable, ainsi qu'une base solide pour tester des évolutions de modèle si
les règles métier évoluent.

\end{document}
